% Part: proof-theory
% Chapter: proof-search
% Section: introduction

\documentclass[../../../include/open-logic-section]{subfiles}

\begin{document}

\olfileid{pt}{ps}{int}

\olsection{تمهيد في البحث عن البرهان}

تبيّن مسلّمات أنظمة~!!{proof} وقواعدها، مثل~\Log{G3c} أو~\Log{N1i}،
أي أشجار المتتاليات أو~!!{formula}s يصح أن تعد~!!{proof}s. فإذا أعطينا
شجرة من هذه المتتاليات أو~!!{formula}s، ومعها عند الحاجة ما يزيد من بيان، كالقواعد المستعملة في كل
خطوة أو وسوم الفروض والقواعد المبرئة لها، أمكننا من تعريف المسلّمات
والقواعد أن نتحقق من كونها~!!{proof} صحيحًا. وأما \emph{العثور} على
!!a{proof} لمتتالية معلومة، أو لـ~!!{formula} من فروض معلومة، فمسألة
أخرى، وهي في الغالب أشق. وتسمى مسألة \emph{البحث عن البرهان}.

وللبحث سبيل بالغ السذاجة. فـ!!{formula}s متتاليات رموز من أبجدية منتهية،
ومن ثم تقبل العد؛ وكذلك المتتاليات وأشجار المتتاليات أو~!!{formula}s،
إذ تمثل هي أيضًا بمتتاليات من~!!{formula}s، وربما زيدت أقواس تدل على
تشعب الشجرة. وبالمسلّمات والقواعد نفحص آليًا هل متتالية الرموز~!!{proof}
صحيح. فيسعنا أن نولد آليًا \emph{جميع}~!!{proof}s الصحيحة الممكنة: نولد
كل متتاليات الرموز من أبجدية تمثيل~!!{proof}s، ثم نمتحن كل مرشح لنعلم
هل هو~!!{proof} صحيح. والبحث الساذج أن نمضي في توليد~!!{proof}s الصحيحة
حتى نجد منها~!!{proof} للمتتالية المطلوبة، أو لـ~!!{formula} المطلوبة
مع فروض مفتوحة لا تجاوز الفروض المعطاة.

والأجدى أن نصطنع الطريقة اليدوية المعتادة في طلب~!!a{proof}: نبدأ
بالمتتالية النهائية، ونختار~!!a{formula}، ثم نعمل قاعدة استدلال «إلى
الخلف» فنولد مقدمة أو مقدمات؛ ولو وجدت لها~!!{proof}s لاتحدت مع آخر
الاستدلال فكانت~!!a{proof} للمتتالية المعطاة. وللبحث عن~!!a{proof} في
نظام استنباط طبيعي طريقة شبيهة بهذا، وإن زادت تعقيدًا.

لكن هذه الطريقة لا ضمان لها؛ فقد تسد عليك السبيل إذا أخطأت اختيار القاعدة
أو ترتيب~!!{formula}s. وهي أيضًا غير متعينة تمامًا، لأن في الموضع الواحد
ربما صلح للاختيار أكثر من~!!{formula} أو للتطبيق إلى الخلف أكثر من قاعدة.
وأما \emph{خوارزمية البحث عن البرهان} فهي إجراء متعين يجد~!!a{proof}
متى كان موجودًا، فهذا معنى ضمانها.

وبعض أنظمة~!!{proof} أوفق من بعض لخوارزميات البحث اليسيرة. ففي أنظمة
المتتاليات يعين~!!{main operator} لـ~!!a{formula} والجانب الذي تقع فيه
القاعدة الواجب إعمالها إلى الخلف. فإذا طلبت، مثلًا،~!!a{proof} للمتتالية
$\OLId{greek-variable}{Gamma} \Sequent \OLId{greek-variable}{Delta},!A\land !B$، وكانت~$!A\land !B$ عن اليمين،
وجب أن تعمل~\RightR{\land} إلى الخلف، فتولد المقدمتان
$\OLId{greek-variable}{Gamma} \Sequent \OLId{greek-variable}{Delta},!A$ و
$\OLId{greek-variable}{Gamma} \Sequent \OLId{greek-variable}{Delta},!B$. أما إذا كان نظامك~\Log{G1c}، فإن وجود
الانقباض يعقّد الأمر لأنه يمنحك خيارات أكثر: يمكنك كذلك تطبيق
\RightR{\Contraction} إلى الخلف وتوليد المقدمة
$\OLId{greek-variable}{Gamma} \Sequent \OLId{greek-variable}{Delta},!A\land !B,!A\land !B$، ثم المقدمة ذات النسخ
الثلاث، وعلى هذا. ويزداد الأمر سوءًا في~\Log{G2c}. فتطبيق~\RightR{\land} إلى
الخلف يقتضي كذلك تخمين~$\OLId{greek-variable}{Gamma}_\OLMathNumeral{1},\OLId{greek-variable}{Gamma}_\OLMathNumeral{2}$ بحيث
$\OLId{greek-variable}{Gamma}=\OLId{greek-variable}{Gamma}_\OLMathNumeral{1}\cup\OLId{greek-variable}{Gamma}_\OLMathNumeral{2}$، وتخمين~$\OLId{greek-variable}{Delta}_\OLMathNumeral{1},\OLId{greek-variable}{Delta}_\OLMathNumeral{2}$ بحيث
$\OLId{greek-variable}{Delta}=\OLId{greek-variable}{Delta}_\OLMathNumeral{1}\cup\OLId{greek-variable}{Delta}_\OLMathNumeral{2}$، ثم تجربة المقدمتين
$\OLId{greek-variable}{Gamma}_\OLMathNumeral{1}\Sequent\OLId{greek-variable}{Delta}_\OLMathNumeral{1},!A$ و$\OLId{greek-variable}{Gamma}_\OLMathNumeral{2}\Sequent\OLId{greek-variable}{Delta}_\OLMathNumeral{2},!B$. وقد يقودك
تخمين خاطئ إلى طريق مسدود لاحقًا، فتضطر إلى الرجوع للبداية واختيار سياقات
مختلفة. وإذا كانت~!!{formula} هي~$!A\lor !B$، فعليك اختيار إحدى صورتي
\RightR{\lor}، فتنتج إحدى المقدمتين
$\OLId{greek-variable}{Gamma}\Sequent\OLId{greek-variable}{Delta},!A$ أو~$\OLId{greek-variable}{Gamma}\Sequent\OLId{greek-variable}{Delta},!B$. ويضطرك الاختيار
الخاطئ مرة أخرى إلى الرجوع.

وأما~\Log{G3c} فسالم من هذه المشكلات؛ إذ لا قواعد بنيوية فيه، واختيار
قاعدة الاستدلال متعين بـ!!{main operator} وبالجانب الذي تقع فيه
!!{formula}. فـ\Log{G3c} أنسب للبحث عن البرهان من~\Log{G1c} أو~\Log{G2c}.
والبحث الناجح ينتج~!!a{proof} في~\Log{G3c}، ثم تنقل~!!{proof} إلى
\Log{G1c} أو~\Log{G2c}، بل إلى~\Log{N1c}. فإذا كانت لـ~\Log{G3c}
خوارزمية بحث، وكانت لنقل~!!{proof}s~\Log{G3c} إلى~!!{proof}s الأنظمة
الأخرى خوارزميات، انحلت بذلك مسألة البحث في تلك الأنظمة.

وبم يعرف ضمان الخوارزمية؟ بأن نثبت أن إخفاقها في العثور على~!!a{proof}
يدل على امتناع~!!{proof}. فقد تعجز خوارزمية رديئة عن برهان موجود. ولا
يعرض هذا للبحث الساذج، لأنه يستقصي جميع~!!{proof}s الممكنة؛ لكن المطلوب
من خوارزمية البحث أن تكون فعالة مقتصدة في الجهد.

وعمدة إثبات الضمان أن نستخرج من كيفية الإخفاق دليلًا على امتناع
!!a{proof} للمتتالية، وذلك بإثبات عدم صلاحها. ففي منطق الرتبة الأولى
الكلاسيكي تكون المتتالية~$\OLId{greek-variable}{Gamma}\Sequent\OLId{greek-variable}{Delta}$ صالحة إذا جعل كل
!!{structure}~!!{formula} واحدة على الأقل من~$\OLId{greek-variable}{Gamma}$ كاذبة، أو جعل
!!{formula} واحدة على الأقل من~$\OLId{greek-variable}{Delta}$ صادقة. والنسق~\Log{G3c}
\emph{سليم}، أي لا~!!{prove}s إلا متتاليات صالحة. وهذا يبرهن بالاستقراء
على~!!{proof}s: فالمسلّمات صالحة، وصلاح مقدمات القاعدة يستلزم صلاح
نتيجتها. ثم نثبت ضمان الخوارزمية بأن نبين أن إخفاق البحث عن~!!a{proof}
للمتتالية~$\OLId{greek-variable}{Gamma}\Sequent\OLId{greek-variable}{Delta}$ يتيح تعريف~!!a{structure} تكون جميع
!!{formula}s~$\OLId{greek-variable}{Gamma}$ فيه صادقة وجميع~!!{formula}s~$\OLId{greek-variable}{Delta}$ كاذبة.
وبذلك يثبت أيضًا أن~\Log{G3c} \emph{كامل}: فإذا كانت
$\OLId{greek-variable}{Gamma}\Sequent\OLId{greek-variable}{Delta}$ صالحة، وجد لها~!!{proof} في~\Log{G3c}. ولما كان
كل بحث فاشل شاهدًا على عدم الصلاح، وجب أن ينجح البحث عن كل متتالية صالحة.

\end{document}
